\documentclass[11pt,a4paper]{article}
\usepackage[margin=1in]{geometry}
\usepackage{amsmath,amssymb,amsthm}
\usepackage{graphicx}
\usepackage{microtype}
\usepackage[hidelinks]{hyperref}

\newtheorem{theorem}{Theorem}
\newtheorem{corollary}[theorem]{Corollary}
\newtheorem{remark}[theorem]{Remark}
\newcommand{\R}{\mathbb R}

\title{Every positive finite Ermakoff limit occurs}
\author{A full answer to the footnote question in arXiv:2305.12873}
\date{17 August 2026}

\begin{document}
\maketitle

\noindent\textbf{Status.}
Candidate full proof, likely valid, pending human review.  For every prescribed
$a\in(0,\infty)$ we construct a continuous strictly increasing function
$g:[1,\infty)\to[1,\infty)$ with $g(1)=1$ such that
\[
 \frac{t g(t)}{g(e^t)}=a
\]
for all sufficiently large $t$.  Thus the set of positive finite Ermakoff
limits is the whole interval $(0,\infty)$.

\section{Source question}

Footnote~3 on source PDF page~5 recalls Ermakoff's test for a continuous
positive increasing function $g:[1,\infty)\to[1,\infty)$ with $g(1)=1$ and
asks for an example in which
\[
 \lim_{t\to\infty}\frac{t g(t)}{g(e^t)}
\]
exists but is different from $0$, $1$, and $\infty$~\cite{MMP}.

\begin{center}
\includegraphics[width=\textwidth]{figures/source_question_crop.png}
\end{center}

\section{Proof intuition}

Two logarithmic changes of variables turn the requested asymptotic ratio
into an additive functional equation.  Write
\[
 h(x)=\log g(e^x),\qquad c=\log a.
\]
For $t=e^x$ the identity $t g(t)/g(e^t)=a$ is equivalent to
\begin{equation}\label{eq:abel}
 h(e^x)=h(x)+x-c.
\end{equation}
This equation can be solved one iterated-exponential interval at a time.
Choose a short increasing seed that has exactly the endpoint jump required
by \eqref{eq:abel}; thereafter the recurrence itself propagates continuity
and strict monotonicity forever.

\section{Construction and proof}

\begin{theorem}\label{thm:main}
For every $a>0$ there is a continuous strictly increasing
$g:[1,\infty)\to[1,\infty)$ with $g(1)=1$ for which
\begin{equation}\label{eq:ratio}
 \frac{t g(t)}{g(e^t)}=a
\end{equation}
for every sufficiently large $t$.
\end{theorem}

\begin{proof}
Fix $a>0$, put $c=\log a$, and choose
\[
 A>\max\{0,c\},\qquad E_0=e^A,
 \qquad m=\frac{A-c}{e^A-A}>0.
\]
Begin with the linear seed
\begin{equation}\label{eq:seed}
 h(x)=mx,\qquad 0\le x\le E_0.
\end{equation}
Set $E_{n+1}=e^{E_n}$.  Successively on each interval
$[E_n,E_{n+1}]$, $n\ge0$, define
\begin{equation}\label{eq:recurrence}
 h(y)=h(\log y)+\log y-c.
\end{equation}
This is a genuine recursive definition: when $y\in[E_n,E_{n+1}]$,
the argument $\log y$ lies in the part of the half-line on which $h$ has
already been defined.  Equivalently, one may iterate logarithms until the
argument first enters $[0,E_0]$ and then apply \eqref{eq:seed}.

We verify the three properties needed below.  At the first junction,
\[
 h(A)+A-c=mA+A-c=mA+m(e^A-A)=mE_0=h(E_0),
\]
so \eqref{eq:recurrence} agrees with the seed.  At every later junction the
two neighboring definitions use the same continuous expression in the
already defined value $h(\log y)$.  Induction therefore proves that $h$ is
continuous on $[0,\infty)$.

The seed is strictly increasing.  If $h$ is increasing on the previously
constructed part and $y_1<y_2$ lie in the next interval, then
\[
 h(y_2)-h(y_1)
 =h(\log y_2)-h(\log y_1)+\log y_2-\log y_1>0.
\]
Continuity at the junctions now proves by induction that $h$ is strictly
increasing everywhere.  Finally, $h\ge0$ on the seed interval, and each
recursive step adds $\log y-c\ge A-c>0$; hence $h$ takes values in
$[0,\infty)$.

For every $x\ge A$, the point $e^x$ lies beyond $E_0$, so the defining
recurrence gives the exact identity
\begin{equation}\label{eq:exact}
 h(e^x)=h(x)+x-c.
\end{equation}
Define
\begin{equation}\label{eq:g}
 g(t)=\exp(h(\log t)),\qquad t\ge1.
\end{equation}
The properties just proved show that $g$ is continuous, strictly increasing,
takes values in $[1,\infty)$, and satisfies $g(1)=1$.  If $t=e^x\ge e^A$,
then \eqref{eq:exact} yields
\[
 \frac{t g(t)}{g(e^t)}
 =\exp\!\bigl(x+h(x)-h(e^x)\bigr)
 =\exp(c)=a.
\]
This proves \eqref{eq:ratio}, in fact as an eventual exact identity rather
than merely a limiting relation.
\end{proof}

\begin{corollary}\label{cor:two}
There is a continuous strictly increasing $g:[1,\infty)\to[1,\infty)$,
$g(1)=1$, for which the Ermakoff limit equals $2$.
\end{corollary}

\begin{proof}
In Theorem~\ref{thm:main}, take $a=2$ and $A=1$.  Then
\[
 m=\frac{1-\log2}{e-1},\qquad
 h(x)=mx\quad(0\le x\le e),
\]
and extend by $h(y)=h(\log y)+\log y-\log2$ for $y\ge e$.
The resulting $g$ from \eqref{eq:g} has
$t g(t)/g(e^t)=2$ for every $t\ge e$.
\end{proof}

\section{Checks, upgrade, novelty, and scope}

The first construction used the explicit value $a=2$, which already answers
the source question.  The deep upgrade in Theorem~\ref{thm:main} shows that
the same mechanism realizes every prescribed $a\in(0,\infty)$.  The choice
of a flexible threshold $A>\max\{0,\log a\}$ is what keeps the seed slope and
every recursive increment positive, even when $a$ is very large.

The script \texttt{code/verify\_construction.py} evaluates the iterated-log
formula, checks the functional equation and eventual ratios for several
values of $a$, checks monotonicity on logarithmic grids, and probes both sides
of the first iterated-exponential junctions.  These numerical checks are not
used in the proof.

A bounded primary-source search used the exact footnote wording,
``Ermakoff's condition'', $t g(t)/g(e^t)$, and the conjugated functional
equation \eqref{eq:abel}.  It found the source question but no primary paper
giving a nontrivial finite example or the arbitrary-$a$ construction.
Novelty confidence is moderate to high; correctness confidence is high.

The theorem answers only the footnote's existence question.  It does not
alter the paper's Poincar\'e inequality theorem, whose hypothesis specifically
requires Ermakoff limit zero.

\begin{thebibliography}{1}
\bibitem{MMP}
J.~Mart\'in and W.~A. Ortiz,
\emph{A note on rearrangement Poincar\'e inequalities and the doubling
condition}, arXiv:2305.12873.
\end{thebibliography}

\end{document}
